\scanpage{51}
Dans le cas où $k$ est de car. $0$, \unsure{celà|cel|ceci|cela} est donc une connexion intégrable absolue sur $\mathscr H^*$. (Manin n'utilise pas ce langage géométrique). Puis, dans le cas du $H^1$, et lorsque $K$ \deleted{est} provient d'un modèle \inserted{propre} lisse, Manin montre de façon assez compliquée que $H^1(X,\Omega^{\cdot}_{X/k})$ (qui est $\subset H^1_{\mathrm{DR}}(K/k)$) est stable par les opérations $\Theta_x$

3.5. En fait, par des arguments très simples directs, \deleted{\unsure{que|qui}} reposant sur la ``formule d'homotopie'' de Cartan
\[
d i_x+i_x d=\Theta_x
\]
on peut montrer, pour un morphisme $f:X\longrightarrow S$ :

(i) Pour toute dérivation $\xi$ de $\underline O_S$ qui se remonte, \deleted{\unsure{sur $X$,|de $X$,|en $X$,}} \inserted{loc. au dessus de $S$,} en une dérivation $\bar\xi$ de $\underline O_X$ [condition tjrs satisfaite si $X$ est \deleted{\unsure{de|est}} \uline{formellt lisse sur} $S$ et \uline{affine} sur $S$] l'effet de $(\bar\xi,\xi)$ opérant sur $\deleted{\unsure{\mathscr H|R}} f_*(\Omega^{\cdot}_{X/S})$, \unsure{est|\deleted{est}} indépendant \uline{mod} \uline{homotopie} \deleted{\unsure{du}} du choix de $\bar\xi$, donc dépend seulement de $\xi$ ; donc en particulier, \inserted{si $f$ formellt lisse et affine} les $\mathscr H^i(f_*(\Omega^{\cdot}_{X/S}))=\mathbb R^i f_*(\Omega^{\cdot}_{X/S})$ sont des Modules sur lesquels les dérivations de $S$ opèrent, -- et on vérifie la loi de linéarité et de crochet, trivialement. Si $S$ lisse sur un $T$, ceci donne pour $f_*(\Omega^{\cdot}_{X/S})$ et \unsure{ts|tous|via} les $R^i f_*(\Omega^{\cdot}_{X/S})$ des conn. intégrables rel/$T$.
% archive sheet number: 50

